\DocumentMetadata{
  pdfversion=2.0,
  pdfstandard=ua-2,
  lang=en-US,
  tagging=on,
  tagging-setup={math/setup=mathml-SE,tabsorder=structure}
}
\documentclass[11pt,letterpaper]{article}
\usepackage[margin=0.68in,includefoot]{geometry}
\usepackage{newcomputermodern}
\usepackage{amsmath,amscd,mathtools}
\usepackage{tikz,adjustbox}
\providecommand{\square}{\mdlgwhtsquare}
\usepackage[hidelinks]{hyperref}
\hypersetup{
  pdftitle={Algebra PhD Exam, May 1995},
  pdfauthor={University of Florida Department of Mathematics},
  pdfsubject={Graduate mathematics examination},
  pdfdisplaydoctitle=true,
  bookmarksopen=true
}
\setcounter{secnumdepth}{0}
\setlength{\parindent}{0pt}
\setlength{\parskip}{0.28em}
\setlength{\emergencystretch}{3em}
\setlength{\leftmargini}{2.15em}
\setlength{\leftmarginii}{2.65em}
\setlength{\leftmarginiii}{3em}
\pagestyle{empty}
\raggedbottom
\allowdisplaybreaks
\NewTaggingSocketPlug{block/list/label}{examproblem}{%
  \tagstructbegin{tag=Lbl}\tagstructend%
  \tagstructbegin{tag=\LBody}%
  \tagstructbegin{tag=\ProblemHeadingTag,title-o={\ProblemHeadingText}}%
  \tagmcbegin{tag=\ProblemHeadingTag,actualtext={\ProblemHeadingText}}%
  #2%
  \tagmcend%
  \tagstructend%
}
\newenvironment{examproblems}[2]{%
  \def\ProblemHeadingTag{#2}%
  \AssignTaggingSocketPlug{block/list/label}{examproblem}%
  \begin{enumerate}%
  \setcounter{enumi}{#1}%
  \renewcommand{\labelenumi}{\textbf{\theenumi.}}%
  \setlength{\topsep}{0.35em}%
  \setlength{\partopsep}{0pt}%
  \setlength{\itemsep}{0.48em}%
  \setlength{\parsep}{0.18em}%
}{\end{enumerate}}
\newenvironment{examsubparts}{%
  \AssignTaggingSocketPlug{block/list/label}{default}%
  \begin{enumerate}%
  \setlength{\topsep}{0.18em}%
  \setlength{\partopsep}{0pt}%
  \setlength{\itemsep}{0.16em}%
  \setlength{\parsep}{0.1em}%
}{\end{enumerate}}
\newenvironment{examinstructions}{%
  \par\begingroup\itshape
}{%
  \par\endgroup\vspace{0.18em}
}
\makeatletter
\renewcommand\subsection{\@startsection{subsection}{2}{0pt}%
  {-1.25ex plus -.25ex minus -.1ex}%
  {0.55ex plus .1ex}%
  {\normalfont\large\bfseries}}
\renewcommand{\@maketitle}{%
  \newpage\null\vskip -1em%
  {\footnotesize\bfseries
    \href{https://www.ufl.edu/}{University of Florida}, \href{https://math.ufl.edu/}{Department of Mathematics}\par}%
  \vskip -0.5em%
  \tagstructbegin{tag=H1,title={Algebra PhD Exam, May 1995}}%
  \tagpdfparaOff%
  \tagmcbegin{tag=H1}%
    {\LARGE\bfseries\href{https://gma.math.ufl.edu/exams/algebra-phd/}{Algebra PhD Exam}, May 1995\par}%
  \tagmcend%
  \tagpdfparaOn%
  \tagstructend%
  \vskip 0.25em%
  \tagmcbegin{artifact}%
    \hrule height 1pt%
  \tagmcend%
  \vskip 1em%
}
\makeatother
\title{Algebra PhD Exam, May 1995}
\author{}
\date{}
\begin{document}
\maketitle
\thispagestyle{empty}
\begin{examinstructions}
Answer seven problems. Write your answers clearly in complete English sentences. You may quote results (within reason) as long as you state them clearly.
\end{examinstructions}
\begin{examproblems}{0}{H2}
\def\ProblemHeadingText{Problem 1.}
\item
Let~\(G\) be a finite group and let~\(A\) be a subgroup of \(\operatorname{Aut}(G)\) such that \(G=[G,A]\). Let~\(N\) be a normal subgroup of~\(G\) such that \([N,A]=1\). Prove that~\(N\) is contained in the center of~\(G\).
\def\ProblemHeadingText{Problem 2.}
\item
This problem has two parts.
\begin{examsubparts}
\item[\textnormal{(a)}]
Give an example of fields \(M\supset L\supset K\) such that \(M/L\) and \(L/K\) are normal extensions, but \(M/K\) is not normal. Justify your answer.
\item[\textnormal{(b)}]
Let \(L\supset K\) be fields such that \([L:K]=8\). Prove that there are \(\alpha,\beta,\gamma\in L\) such that \(L=K(\alpha,\beta,\gamma)\).
\end{examsubparts}
\def\ProblemHeadingText{Problem 3.}
\item
Let~\(A\) be a semi-simple commutative ring with 1. Prove that if~\(A\) is Artinian then~\(A\) is Noetherian. Give a counterexample to show that the converse is not true.
\def\ProblemHeadingText{Problem 4.}
\item
This problem has two parts.
\begin{examsubparts}
\item[\textnormal{(a)}]
Give an example of a free module~\(F\) over a ring with 1 which has a submodule~\(M\) which is not free.
\item[\textnormal{(b)}]
Let~\(A\) be a commutative ring with 1 and let~\(P\) and~\(Q\) be projective~\(A\)-modules. Prove that \(P\otimes Q\) is projective.
\end{examsubparts}
\def\ProblemHeadingText{Problem 5.}
\item
Let~\(A\) be a commutative ring with 1 and let \(S\not\ni 0\) be a multiplicative subset of~\(A\). Let~\(P\) be a maximal element in the set \(\{I:I\text{ is an ideal of }A\text{ and }I\cap S=\varnothing\}\). Prove that~\(P\) is a prime ideal.
\def\ProblemHeadingText{Problem 6.}
\item
Let~\(A\) be a commutative ring.
\begin{examsubparts}
\item[\textnormal{(a)}]
Prove that if \(I_1\) and \(I_2\) are ideals in~\(A\) and~\(P\) is a prime ideal in~\(A\) such that \(I_1\cap I_2=P\), then \(I_1=P\) or \(I_2=P\).
\item[\textnormal{(b)}]
Is the above statement still true if the word “prime” is replaced by “primary”? Justify your answer.
\end{examsubparts}
\def\ProblemHeadingText{Problem 7.}
\item
Let~\(n\) and~\(m\) be positive integers. Find~\(r\) such that \(\mathbb{Z}_n\otimes_{\mathbb{Z}}\mathbb{Z}_m\cong\mathbb{Z}_r\). Justify your answer.
\def\ProblemHeadingText{Problem 8.}
\item
Let~\(R\) be a commutative Noetherian ring with 1. Prove that \(R[[X]]\) is Noetherian.
\def\ProblemHeadingText{Problem 9.}
\item
Let~\(R\) be a commutative ring with 1 and let \(S\subset R\) be a multiplicative set which contains no zero divisors. Prove that the prime ideals of the ring \(S^{-1}R\) are in one-to-one correspondence with the prime ideals~\(P\) of~\(R\) such that \(P\cap S=\varnothing\).
\def\ProblemHeadingText{Problem 10.}
\item
Let~\(G\) be an abelian group which is injective. Prove that~\(G\) is divisible.
\def\ProblemHeadingText{Problem 11.}
\item
Let \(f(X)=2X^5-10X+5\in\mathbb{Q}[X]\) and let~\(L\) be a splitting field for~\(f\). Compute the Galois group \(\operatorname{Gal}(L/\mathbb{Q})\).
\end{examproblems}
\end{document}
